\documentclass[12pt, letterpaper]{article}

% --- Paquetes ---
\usepackage[utf8]{inputenc}
\usepackage[spanish]{babel}
\usepackage[margin=2.5cm]{geometry}
\usepackage{amsmath}
\usepackage{graphicx}
\usepackage{booktabs}
\usepackage{array}
\usepackage{fancyhdr}
\usepackage{titlesec}
\usepackage{enumitem}
\usepackage{parskip}

% --- Encabezado y pie de página ---
\pagestyle{fancy}
\fancyhf{}
\fancyhead[L]{\small CC3182 -- Visión por Computadora}
\fancyhead[R]{\small Laboratorio 7}
\fancyfoot[C]{\thepage}
\renewcommand{\headrulewidth}{0.4pt}

% --- Formato de secciones: solo negrita, sin colores ---
\titleformat{\section}{\large\bfseries}{}{0em}{}[\titlerule]
\titleformat{\subsection}{\normalsize\bfseries}{}{0em}{}

% --- Inicio del documento ---
\begin{document}

% ===================== PORTADA =====================
\begin{titlepage}
    \centering
    \vspace*{1cm}

    {\large \textbf{Universidad del Valle de Guatemala}\\
    Facultad de Ingeniería\\
    CC3182 -- Visión por Computadora}

    \vspace{2cm}

    {\Huge \textbf{Laboratorio 7}}\\[0.5cm]
    {\Large Transfer Learning y Entrenamiento de Modelos}

    \vspace{2cm}

    \begin{tabular}{ll}
        \textbf{Integrante 1:} & Andy Fuentes \\[4pt]
        \textbf{Integrante 2:} & Christian Echeverría \\[4pt]
        \textbf{Integrante 3:} & Diederich Solís \\
    \end{tabular}

    \vspace{1.5cm}

    \textbf{Fecha:26/03/2026} \
    
    \textbf{Repositorio:https://github.com/DiederichSolis/LAB7vs.git} \

    \vfill

    {\small Dataset utilizado: \textit{Chest X-Ray Images (Pneumonia)} -- Kaggle}
\end{titlepage}

% ===================== TASK 1 =====================
\section*{\LARGE Task 1: Análisis de Decisiones de Entrenamiento}

\noindent

% ---- PREGUNTA 1.1 ----
\section{Pregunta 1.1 -- Data Augmentation}

\subsection{Pregunta 1: ¿Qué es el Data Augmentation y por qué no son datos falsos?}

Para explicarle al médico coordinador qué es el Data Augmentation, usaríamos la siguiente analogía:

\begin{quote}
\itshape
Imagine que usted tiene 800 fotografías de radiografías y necesita que un residente de medicina
aprenda a identificar neumonía. ¿Qué haría si solo tuviera esas 800 fotos? Probablemente las
mostraría desde distintos ángulos, con diferente iluminación, recortadas de distintas formas,
para que el residente aprenda a reconocer la enfermedad sin importar cómo venga presentada la imagen.
Eso es exactamente el Data Augmentation: tomamos las mismas radiografías reales y las mostramos
de maneras ligeramente distintas al modelo, como si el mismo paciente hubiera sido fotografiado
con el equipo un poco inclinado, o con el brillo ajustado.
\end{quote}

Las imágenes generadas no son datos falsos porque \textbf{no se inventa información diagnóstica nueva}.
Una radiografía rotada 10 grados sigue mostrando el mismo pulmón, las mismas opacidades, el mismo
diagnóstico. Lo único que cambia es la perspectiva, y eso es exactamente lo que queremos: que el
modelo aprenda a reconocer neumonía independientemente de cómo esté orientada o iluminada la imagen.
La patología es la misma; solo varía la forma en que la vemos.

\subsection{Pregunta 2: Transformaciones válidas e inválidas para radiografías de tórax}

A continuación se proponen tres transformaciones adecuadas para este dominio médico, y una que
\textbf{no} debería aplicarse:

\medskip

\textbf{Transformaciones válidas:}

\begin{enumerate}[leftmargin=1.5cm]
    \item \textbf{Rotación pequeña ($\pm$10°):} En la práctica clínica, los pacientes no siempre
    están perfectamente alineados cuando se toma la radiografía. Una ligera rotación simula esa
    variación real. Se mantiene dentro de un rango pequeño para no distorsionar la anatomía
    de referencia (como la posición del corazón o el diafragma).

    \item \textbf{Ajuste de brillo y contraste:} Los equipos de rayos X en diferentes clínicas
    pueden producir imágenes con exposiciones distintas. Variar el brillo y el contraste entrena
    al modelo para ser robusto ante diferencias de calibración entre equipos, algo muy relevante
    para clínicas rurales que pueden no tener el mismo equipo estandarizado.

    \item \textbf{Zoom aleatorio suave (hasta 10--15\%):} Simula variaciones en la distancia
    del paciente al detector o diferencias en el campo de visión según el técnico. Permite que
    el modelo aprenda a identificar opacidades sin depender de que el pulmón ocupe exactamente
    el mismo espacio en la imagen.
\end{enumerate}

\medskip

\textbf{Transformación que NO debería aplicarse:}

\textbf{Volteo horizontal (flip):} Esta transformación sería problemática porque el corazón
está ubicado en el lado izquierdo del tórax. Si se aplica un volteo horizontal, el modelo podría
ver imágenes donde el corazón aparece en el lado derecho, lo cual correspondería a una condición
médica real llamada \textit{dextrocardia}. Esto podría confundir al modelo y, peor aún, podría
enseñarle patrones anatómicos incorrectos que luego aplique en producción, comprometiendo
directamente la integridad diagnóstica.

\subsection{Pregunta 3: ¿Es el Data Augmentation suficiente por sí solo?}

No, el Data Augmentation solo no es suficiente para garantizar que el modelo generalice bien,
aunque es una herramienta muy importante. Existen otras variables críticas que también influyen:

\begin{itemize}[leftmargin=1.5cm]
    \item \textbf{Calidad del etiquetado:} Si las 800 radiografías tienen errores en las etiquetas
    (por ejemplo, casos de neumonía marcados como normales), el augmentation va a multiplicar esos
    errores. Más variaciones de datos mal etiquetados no ayudan.

    \item \textbf{Regularización:} El augmentation ayuda a reducir el overfitting, pero técnicas
    como Dropout, L2 o Early Stopping complementan ese efecto desde la perspectiva del modelo
    mismo. Sin ellas, un modelo puede seguir memorizando incluso con datos aumentados.

    \item \textbf{Arquitectura del modelo y Transfer Learning:} Con solo 800 imágenes, la
    arquitectura importa mucho. Usar un modelo pre-entrenado con Transfer Learning puede suplir
    parte de la escasez de datos al ya traer ``conocimiento visual'' de millones de imágenes.

    \item \textbf{Representatividad del dataset original:} Si las 800 radiografías provienen de
    un solo hospital con un solo equipo de rayos X, el augmentation no va a simular la variabilidad
    real de distintos equipos y poblaciones. La diversidad base del dataset también importa.
\end{itemize}

En resumen, el Data Augmentation es una pieza del rompecabezas, pero no el rompecabezas completo.

% ---- PREGUNTA 1.2 ----
\section{Pregunta 1.2 -- Overfitting en las Curvas de Loss}

Los siguientes datos fueron registrados durante el entrenamiento del modelo de MediScan Guatemala:

\begin{table}[h!]
\centering
\begin{tabular}{ccc}
\toprule
\textbf{Época} & \textbf{Loss de Entrenamiento} & \textbf{Loss de Validación} \\
\midrule
5  & 0.52 & 0.55 \\
10 & 0.31 & 0.38 \\
15 & 0.18 & 0.42 \\
20 & 0.09 & 0.61 \\
25 & 0.04 & 0.89 \\
\bottomrule
\end{tabular}
\caption{Curvas de loss durante 25 épocas de entrenamiento.}
\end{table}

\subsection{Pregunta 4: ¿A partir de qué época comienza el sobreajuste?}

El sobreajuste comienza a evidenciarse claramente a partir de la \textbf{época 15}. La señal más
clara es la divergencia entre las dos curvas de loss:

\begin{itemize}[leftmargin=1.5cm]
    \item En la época 10, ambas métricas siguen bajando juntas (0.31 vs 0.38), lo cual es
    comportamiento normal de un modelo que todavía está aprendiendo de forma generalizable.

    \item En la época 15, el loss de entrenamiento baja a 0.18, pero el de validación sube a 0.42.
    Aquí el modelo ya está mejorando en los datos que conoce, pero empeorando en los que no ha visto.
    Eso es la definición de overfitting.

    \item Para la época 25, la brecha es dramática: 0.04 de entrenamiento vs 0.89 de validación.
    El modelo memorizó casi a la perfección los 800 ejemplos, pero no aprendió nada útil para
    radiografías nuevas.
\end{itemize}

\subsection{Pregunta 5: Estrategias de regularización para prevenir el overfitting}

\begin{enumerate}[leftmargin=1.5cm]

    \item \textbf{Dropout:} Esta técnica consiste en ``apagar'' aleatoriamente un porcentaje de
    neuronas durante cada paso de entrenamiento. El fenómeno matemático que mitiga es la
    \textit{co-adaptación} entre neuronas, es decir, evita que ciertas neuronas dependan tanto
    de otras que se especialicen en memorizar ejemplos específicos. Con Dropout, cada neurona
    tiene que aprender a ser útil por sí sola. En las curvas de la tabla, esperaríamos ver que
    el loss de validación deje de subir a partir de la época 15 y se mantenga más estable, aunque
    el de entrenamiento quizás baje un poco más lento.

    \item \textbf{Early Stopping:} Aunque técnicamente no es una penalización matemática sobre
    los pesos, sí es una forma de regularización muy efectiva. Consiste en detener el entrenamiento
    cuando el loss de validación deja de mejorar (o comienza a empeorar) por un número determinado
    de épocas (``paciencia''). Lo que mitiga es el \textit{sobre-entrenamiento temporal}: el modelo
    aprende cosas útiles hasta cierto punto, y si lo dejamos continuar solo sigue memorizando
    ruido. En la tabla, habríamos parado alrededor de la época 12--13, donde la validación
    todavía estaba mejorando. Eso habría dado un modelo mucho más generalizable.

\end{enumerate}

\subsection{Pregunta 6: ¿Por qué es peligroso desplegar este modelo en producción?}

Desde una perspectiva médica, este patrón de sobreajuste es especialmente peligroso porque el
modelo aprendió a reconocer las \textbf{800 radiografías específicas} que le mostramos, no la
enfermedad en sí misma. Cuando llegue una radiografía de un paciente nuevo --tomada con un equipo
diferente, en una posición ligeramente distinta, de una persona de otra edad o complexión-- el
modelo fallará.

El problema no es solo técnico; es clínico. Un falso negativo en este contexto significa decirle
a un médico que el paciente está bien cuando en realidad tiene neumonía. En clínicas rurales donde
el sistema de MediScan puede ser el único apoyo diagnóstico disponible, ese error puede resultar
en que un paciente con neumonía no reciba tratamiento oportuno, con consecuencias potencialmente
graves o fatales.

Además, un modelo con overfitting extremo tiende a ser impredecible: puede acertar en algunos
casos pero fallar estrepitosamente en otros sin ninguna señal de advertencia. Esa inconsistencia
es incompatible con un entorno clínico donde la confiabilidad es fundamental.

% ---- PREGUNTA 1.3 ----
\section{Pregunta 1.3 -- Métricas y Clases Desbalanceadas}

El dataset de prueba contiene \textbf{700 radiografías normales} y \textbf{150 con neumonía},
para un total de 850 imágenes.

\subsection{Pregunta 7: Accuracy del modelo naive y qué revela sobre el 94\%}

Un modelo ``naive'' que siempre predice ``Normal'' sin analizar nada acertaría en todas las 700
imágenes normales y fallaría en las 150 de neumonía. Su accuracy sería:

\[
\text{Accuracy}_{\text{naive}} = \frac{700}{850} \approx \mathbf{82.35\%}
\]

Esto revela algo importante: si un modelo que \textbf{ni siquiera analiza las imágenes} alcanza
un 82\%, entonces el 94\% reportado no es tan impresionante. La diferencia entre ambos es de apenas
12 puntos porcentuales. Peor aún, no sabemos si ese porcentaje adicional de mejora viene de detectar
correctamente casos de neumonía, o simplemente de clasificar bien un poco más de imágenes normales.
El accuracy como métrica única nos oculta completamente si el modelo está detectando la enfermedad
que nos importa.

\subsection{Pregunta 8: ¿Por qué F1-Score y Sensibilidad son más informativas?}

En problemas médicos con clases desbalanceadas, el accuracy puede ser engañoso porque mide
el desempeño global sin distinguir entre clases. Las métricas más informativas son:

\begin{itemize}[leftmargin=1.5cm]
    \item \textbf{Sensibilidad (Recall):} Mide qué proporción de los casos de neumonía reales
    el modelo logró identificar correctamente. Su fórmula es:
    \[
    \text{Recall} = \frac{\text{Verdaderos Positivos (VP)}}{\text{Verdaderos Positivos (VP)} + \text{Falsos Negativos (FN)}}
    \]
    Clínicamente, una baja sensibilidad significa muchos falsos negativos, es decir, pacientes con
    neumonía a quienes se les dice que están sanos. Ese es el error más peligroso en este contexto,
    porque podría llevar a no tratar una enfermedad que sí existe.

    \item \textbf{F1-Score:} Es el promedio armónico entre Precisión y Recall:
    \[
    F1 = 2 \cdot \frac{\text{Precisión} \times \text{Recall}}{\text{Precisión} + \text{Recall}}
    \]
    El F1 es útil porque obliga a que el modelo sea bueno en ambas cosas: no solo detectar todos
    los casos de neumonía (recall alto), sino también no alarmar al médico con falsos positivos
    innecesarios (precisión alta). En clínicas rurales con recursos limitados, un exceso de
    alarmas falsas puede generar desconfianza en el sistema y sobrecargar al médico.
\end{itemize}

En resumen, el accuracy es una métrica de conveniencia; el Recall y el F1 son métricas de
\textit{responsabilidad clínica}.

\subsection{Pregunta 9: Respuesta honesta al inversionista}

\begin{quote}
\itshape
``El 94\% de accuracy de nuestro modelo es alentador, pero queremos ser transparentes: dado que
la mayoría de las imágenes de prueba son de pacientes sanos, incluso un sistema que nunca detecte
neumonía podría alcanzar el 82\% de accuracy. Lo que realmente importa para un sistema médico es
qué tan bien identifica los casos de enfermedad cuando sí existen, y para eso utilizamos métricas
como la Sensibilidad y el F1-Score, que miden justo eso. Nuestro siguiente paso es optimizar el
modelo específicamente en esas métricas, y someterlo a validación con radiólogos antes de cualquier
despliegue. El objetivo no es impresionar con un número; es construir un sistema en el que los
médicos de las clínicas rurales puedan confiar.''
\end{quote}

% ===================== TASK 2 =====================
\newpage
\section*{\LARGE Task 2: Transfer Learning y Fine-Tuning}

\noindent

% ---- PREGUNTA 2.1 ----
\section{Pregunta 2.1 -- Transfer Learning desde ImageNet}

\subsection{Pregunta 1: ¿Está de acuerdo con el desarrollador?}

No, no estamos de acuerdo con el desarrollador. Su argumento parte de una idea intuitiva pero
incorrecta: que porque ImageNet tiene fotos de perros y gatos, lo que la red aprendió no sirve
para radiografías. La realidad es que las capas convolucionales de una CNN no aprenden ``perros''
o ``gatos'' directamente; aprenden \textbf{jerarquías de características visuales} que son
universales.

Las primeras capas de una red pre-entrenada en ImageNet aprenden detectores de bordes,
gradientes, texturas y patrones de frecuencia. Estas son características fundamentales de
\textit{cualquier} imagen, incluyendo radiografías. Un borde es un borde sin importar si
aparece en la oreja de un gato o en el contorno de un pulmón. Las capas intermedias aprenden
combinaciones más complejas como formas geométricas, patrones repetitivos y relaciones
espaciales, que también son relevantes para detectar opacidades o consolidaciones pulmonares.

Solo las capas finales (las más profundas) son las que se especializan en clasificar objetos
específicos de ImageNet. Pero esas son precisamente las capas que \textbf{reemplazamos} cuando
hacemos Transfer Learning. Descartamos el cabezal de clasificación de 1000 clases de ImageNet
y lo sustituimos por uno adaptado a nuestro problema binario (Normal vs.\ Neumonía).

El hecho de que las radiografías sean imágenes en escala de grises tampoco invalida la
transferencia: las redes pre-entrenadas en ImageNet esperan 3 canales (RGB), pero simplemente
replicamos el canal de grises tres veces. La información de bajo nivel que la red ya aprendió
(contrastes, bordes, texturas) sigue siendo perfectamente aplicable. Múltiples estudios en
imagen médica han demostrado que modelos pre-entrenados en ImageNet superan consistentemente
a los entrenados desde cero cuando se dispone de datasets pequeños, precisamente porque esa
base de conocimiento visual genérico actúa como un ``punto de partida informado'' que reduce
drásticamente la cantidad de datos necesarios para converger.

\subsection{Pregunta 2: Costo real de entrenar desde cero para MediScan Guatemala}

Más allá del costo computacional (que ya de por sí es significativo: entrenar una ResNet50
desde cero requiere cientos de miles de imágenes y muchas horas de GPU), hay al menos dos
dimensiones críticas que hacen inviable esta decisión para una startup como MediScan Guatemala:

\begin{enumerate}[leftmargin=1.5cm]
    \item \textbf{Tiempo al mercado (time-to-market):} Una startup médica compite contra el
    reloj. Cada mes adicional de desarrollo sin un producto funcional es un mes sin ingresos,
    sin validación clínica y sin retroalimentación de usuarios reales. Con Transfer Learning,
    se puede tener un modelo funcional en días o semanas. Entrenando desde cero con solo 800
    imágenes, no solo tomaría mucho más tiempo experimentar con arquitecturas, sino que el
    resultado probablemente sería inferior. Mientras MediScan está intentando que su modelo
    desde cero aprenda qué es un borde, un competidor con Transfer Learning ya está en fase
    de pruebas clínicas.

    \item \textbf{Riesgo técnico y de rendimiento:} Con 800 imágenes, un modelo entrenado desde
    cero tiene una probabilidad altísima de sobreajustar, exactamente como vimos en el Task~1
    con las curvas de loss. Los pesos inicializados aleatoriamente necesitan grandes cantidades
    de datos para converger a representaciones útiles. En cambio, los pesos pre-entrenados ya
    codifican millones de patrones visuales, lo que actúa como una forma implícita de
    regularización. El riesgo de obtener un modelo clínicamente inútil es mucho mayor
    entrenando desde cero, y en un contexto médico ese riesgo se traduce directamente en
    diagnósticos fallidos.
\end{enumerate}

Una tercera dimensión relevante es el \textbf{talento humano}: diseñar y depurar una
arquitectura desde cero requiere experiencia profunda en deep learning. Con Transfer Learning,
un equipo con conocimientos sólidos pero no necesariamente expertos en diseño de arquitecturas
puede lograr resultados competitivos usando modelos probados por la comunidad investigadora.

\subsection{Pregunta 3: ¿Existe un escenario legítimo para entrenar desde cero?}

Sí, existe un escenario legítimo: cuando se dispone de un \textbf{dataset masivo} (decenas o
cientos de miles de imágenes) de un dominio que es \textbf{fundamentalmente diferente} a las
imágenes naturales de ImageNet, y cuando además se tiene acceso a los recursos computacionales
necesarios.

Un ejemplo concreto sería si MediScan Guatemala, después de varios años de operación, lograra
recopilar más de 100,000 radiografías etiquetadas por radiólogos certificados, y además quisiera
construir un modelo especializado que no solo detecte neumonía sino múltiples patologías
pulmonares con alta granularidad (tuberculosis, derrame pleural, cardiomegalia, nódulos, etc.).
En ese caso, entrenar una arquitectura diseñada específicamente para radiografías ---quizás
con un solo canal de entrada, bloques convolucionales optimizados para el rango dinámico de
rayos X, y un esquema de atención adaptado a la anatomía torácica--- podría superar al Transfer
Learning porque la red aprendería representaciones internas \textit{nativamente} optimizadas
para ese dominio, sin arrastrar sesgos de ImageNet (como la tendencia de las capas tempranas
a sobre-representar texturas de color que no existen en radiografías).

Sin embargo, para el estado actual de MediScan Guatemala ---800 imágenes, recursos limitados,
necesidad de resultados rápidos--- entrenar desde cero sería una decisión técnicamente
injustificable.

% ---- PREGUNTA 2.2 ----
\section{Pregunta 2.2 -- Estrategia de Fine-Tuning}

\subsection{Pregunta 4: ¿Opción A o B para el dataset de MediScan Guatemala?}

Recomendamos la \textbf{Opción A: congelar toda la red base} y entrenar únicamente el cabezal
de clasificación, al menos como punto de partida.

La lógica detrás de esta recomendación se basa en la interacción entre el tamaño del dataset
y la cantidad de parámetros entrenables:

\begin{itemize}[leftmargin=1.5cm]
    \item Con solo \textbf{800 imágenes}, la cantidad de datos es extremadamente limitada.
    Si descongelamos toda la red (Opción~B), estamos exponiendo millones de parámetros
    (ResNet50 tiene aproximadamente 23 millones) a un proceso de optimización con muy pocos
    ejemplos. Esto es una receta directa para el sobreajuste: la red tiene mucha más capacidad
    de la que los datos pueden supervisar.

    \item Al congelar las capas base, estamos aprovechando las representaciones visuales ya
    aprendidas de ImageNet (bordes, texturas, formas) sin arriesgarnos a destruirlas. Solo
    entrenamos el cabezal final, que tiene muchos menos parámetros, lo que reduce drásticamente
    el riesgo de overfitting.

    \item La tasa de aprendizaje alta (1e-3) de la Opción~A es razonable \textit{solo para el
    cabezal}, porque sus pesos se inicializan desde cero y necesitan ajustes grandes para
    converger. Las capas congeladas no se ven afectadas por esta tasa porque sus gradientes
    están desactivados.
\end{itemize}

Dicho esto, la Opción~A pura puede no ser suficiente si el dominio de radiografías es
demasiado distinto de ImageNet. Una estrategia híbrida razonable sería: primero entrenar con
Opción~A hasta que el cabezal converja, y luego descongelar las últimas 2--3 capas de la red
base con una tasa muy baja (1e-5) para hacer un ajuste fino controlado. Esto combina la
seguridad de A con la adaptabilidad parcial de B.

\subsection{Pregunta 5: Riesgo de la Opción B con tasa de aprendizaje alta}

Si aplicamos la Opción~B con una tasa de aprendizaje alta (1e-3) en \textbf{toda} la red, el
riesgo principal es el \textbf{catastrophic forgetting} (olvido catastrófico). Este fenómeno
ocurre cuando los gradientes grandes generados por una tasa de aprendizaje alta modifican
drásticamente los pesos pre-entrenados, destruyendo las representaciones útiles que la red ya
había aprendido de ImageNet.

En términos concretos, lo que sucedería es lo siguiente: las capas tempranas que ya sabían
detectar bordes, texturas y patrones básicos serían sobrescritas con pesos ajustados
exclusivamente a las 800 radiografías de entrenamiento. El resultado es que la red ``olvida''
todo su conocimiento previo y se convierte, efectivamente, en un modelo entrenado desde cero
pero peor, porque partió de pesos que fueron \textit{desestabilizados} en lugar de
inicializados de forma controlada.

Durante el entrenamiento, esto se evidenciaría de varias formas:

\begin{itemize}[leftmargin=1.5cm]
    \item El loss de entrenamiento podría \textbf{aumentar abruptamente} en las primeras épocas,
    porque los pesos pre-entrenados se están destruyendo y la red pierde temporalmente su
    capacidad de extraer características útiles.

    \item Después de ese pico, el modelo podría empezar a bajar el loss de entrenamiento
    rápidamente, pero el loss de validación no lo seguiría, mostrando un patrón de
    \textbf{overfitting severo} similar (o peor) al que vimos en la tabla del Task~1.

    \item La métrica de accuracy en validación podría ser \textbf{inestable}, con oscilaciones
    grandes entre épocas, lo cual es un signo clásico de que la tasa de aprendizaje es
    demasiado alta para las capas que se están ajustando.
\end{itemize}

Por eso la Opción~B original propone una tasa de 1e-5: lo suficientemente baja para hacer
ajustes finos sin destruir lo aprendido. Pero con 800 imágenes, incluso esa tasa es riesgosa
si se aplica a toda la red.

\subsection{Pregunta 6: ¿Cambiaría la recomendación con 50,000 radiografías?}

Sí, con 50,000 radiografías etiquetadas la recomendación cambiaría significativamente. La
estrategia de Fine-Tuning evolucionaría de la siguiente manera:

\begin{enumerate}[leftmargin=1.5cm]
    \item \textbf{Fase 1 -- Cabezal (igual que antes):} Seguiríamos empezando por congelar toda
    la red base y entrenar el cabezal de clasificación con una tasa de aprendizaje alta (1e-3).
    Esto establece una buena línea base y permite que el cabezal aprenda a mapear las
    representaciones existentes al problema de neumonía.

    \item \textbf{Fase 2 -- Descongelamiento progresivo:} Con 50,000 imágenes, ahora sí
    tendríamos suficientes datos para supervisar el ajuste de capas más profundas sin
    sobreajustar. Descongelaríamos la red por bloques, empezando por las capas más profundas
    (cercanas al cabezal) y avanzando hacia las más tempranas. Cada bloque descongelado se
    entrenaría con una tasa de aprendizaje decreciente (discriminative learning rates):
    1e-4 para las últimas capas, 1e-5 para las intermedias, 1e-6 para las tempranas.

    \item \textbf{Fase 3 -- Fine-Tuning completo:} Con esta cantidad de datos, eventualmente
    podríamos descongelar toda la red con una tasa global muy baja (1e-5 o menor) para un
    ajuste final. Esto permitiría que incluso las capas tempranas se adapten ligeramente al
    dominio de rayos X, aprendiendo detectores de bordes y texturas optimizados para el rango
    dinámico y las características específicas de las radiografías.
\end{enumerate}

La razón del cambio es simple: con más datos, el riesgo de overfitting disminuye
proporcionalmente, y el beneficio de adaptar capas profundas al dominio específico aumenta.
Con 50,000 imágenes, la red tendría suficientes ejemplos para aprender representaciones
médicas especializadas sin destruir la estructura general aprendida de ImageNet. Es un
balance: más datos permiten más libertad para ajustar sin perder generalización.

% ---- PREGUNTA 2.3 ----
\section{Pregunta 2.3 -- Reutilización del Modelo para Fracturas}

\subsection{Pregunta 7: Viabilidad de reutilizar el modelo de neumonía para fracturas de muñeca}

La reutilización es \textbf{parcialmente viable}, pero no directa. Para entender qué partes
del modelo podemos aprovechar, debemos analizar qué ha aprendido cada sección de la red:

\begin{itemize}[leftmargin=1.5cm]
    \item \textbf{Capas tempranas (aprovechables):} Los detectores de bordes, gradientes,
    texturas y patrones de bajo nivel que la red aprendió (primero de ImageNet y luego
    refinados para radiografías de tórax) son perfectamente aplicables a radiografías de
    muñeca. Un borde óseo, una línea de fractura, un cambio de densidad ---todo esto depende
    de las mismas primitivas visuales que las capas tempranas ya codifican.

    \item \textbf{Capas intermedias (parcialmente aprovechables):} Estas capas han aprendido
    patrones espaciales y combinaciones de características que son más específicas al tórax:
    la forma de los pulmones, la posición del corazón, la distribución de opacidades. Algunas
    de estas representaciones podrían ser útiles (como la capacidad de detectar cambios de
    densidad en tejido óseo), pero otras serían irrelevantes para la anatomía de la muñeca.

    \item \textbf{Capas finales y cabezal (no aprovechables):} El cabezal de clasificación
    está entrenado específicamente para distinguir ``Normal'' de ``Neumonía'' en tórax.
    Esto no tiene ninguna utilidad para detectar fracturas en muñeca. Además, la anatomía
    que las capas más profundas aprendieron a reconocer (pulmones, diafragma) es completamente
    diferente a la de la muñeca (radio, cúbito, huesos del carpo).
\end{itemize}

En resumen: el modelo no puede ``reutilizarse directamente'' como sugiere el hospital, pero sí
puede servir como un excelente punto de partida para Transfer Learning hacia el nuevo dominio.

\subsection{Pregunta 8: Plan de acción en tres pasos}

\begin{enumerate}[leftmargin=1.5cm]
    \item \textbf{Paso 1 -- Reemplazar el cabezal de clasificación:} Eliminar la capa final
    del modelo de neumonía y sustituirla por un nuevo cabezal adaptado al problema de fracturas.
    Si el nuevo problema es binario (fractura vs.\ normal), el cabezal tendría una estructura
    similar (capas densas con activación sigmoide), pero con pesos inicializados desde cero.
    Si se requiere clasificación más granular (tipo de fractura, ubicación), el cabezal debería
    tener más neuronas de salida.

    \item \textbf{Paso 2 -- Congelar capas tempranas e intermedias, entrenar el cabezal:}
    Congelar todas las capas convolucionales de la red base y entrenar únicamente el nuevo
    cabezal con una tasa de aprendizaje de 1e-3. Esto permite que el modelo aprenda a mapear
    las representaciones visuales existentes al nuevo problema sin destruir los detectores
    de bajo nivel. Monitorear el loss de validación con Early Stopping (paciencia de 5--7
    épocas).

    \item \textbf{Paso 3 -- Descongelamiento selectivo de las últimas capas:} Una vez que el
    cabezal haya convergido, descongelar las últimas 2--3 capas convolucionales (las más
    profundas) y re-entrenar con una tasa de aprendizaje baja (1e-5). Esto permite que la
    red adapte sus representaciones de alto nivel a la anatomía de la muñeca, que es
    estructuralmente diferente al tórax. Las capas tempranas permanecen congeladas porque
    sus detectores de bordes y texturas son igualmente válidos para ambos dominios.
\end{enumerate}

\subsection{Pregunta 9: Riesgo ético/clínico y validación mínima}

Reutilizar un modelo entrenado en un dominio para otro dominio médico sin validación rigurosa
implica varios riesgos serios:

\begin{itemize}[leftmargin=1.5cm]
    \item \textbf{Sesgo de dominio oculto:} El modelo puede arrastrar sesgos aprendidos del
    dominio original. Por ejemplo, si el modelo de neumonía aprendió a asociar ciertas
    densidades con ``Normal'', podría aplicar ese mismo criterio erróneamente a una radiografía
    de muñeca donde una densidad similar indica una fractura por estrés. Estos sesgos son
    difíciles de detectar sin validación formal.

    \item \textbf{Falsa confianza:} Si el modelo muestra métricas de accuracy razonables en
    pruebas iniciales, el equipo podría asumir que funciona bien, cuando en realidad puede
    estar fallando en casos atípicos o en tipos específicos de fractura que no estaban bien
    representados en los datos de entrenamiento.

    \item \textbf{Responsabilidad clínica:} Un diagnóstico erróneo de fractura puede llevar a
    que un paciente no reciba inmovilización, lo que puede resultar en consolidación incorrecta,
    daño permanente o necesidad de cirugía posterior. En el otro extremo, un falso positivo
    podría generar tratamientos e inmovilizaciones innecesarias.
\end{itemize}

El proceso de validación mínimo antes de cualquier despliegue debería incluir:

\begin{enumerate}[leftmargin=1.5cm]
    \item \textbf{Validación con dataset independiente:} Evaluar el modelo con un conjunto de
    radiografías de muñeca que \textbf{no} se usaron durante el entrenamiento ni el Fine-Tuning,
    idealmente provenientes de diferentes hospitales y equipos de rayos X para evaluar
    generalización.

    \item \textbf{Revisión por especialistas:} Someter los resultados a revisión por
    traumatólogos y radiólogos que evalúen no solo las métricas globales, sino los casos
    específicos donde el modelo falla. Esto permite identificar patrones de error
    clínicamente peligrosos que las métricas agregadas no capturan.

    \item \textbf{Fase piloto controlada:} Desplegar el modelo en modo de ``segunda opinión''
    (nunca como único diagnóstico) durante un período de prueba, donde cada predicción es
    verificada por un médico. Solo después de demostrar consistencia y seguridad en esta fase
    se podría considerar un despliegue más amplio.

    \item \textbf{Métricas específicas por tipo de fractura:} No basta con reportar accuracy
    global. Se deben calcular sensibilidad y especificidad desglosadas por tipo de fractura
    (Colles, Smith, escafoides, etc.), porque el modelo podría funcionar bien para fracturas
    obvias pero fallar en fracturas sutiles que son las más difíciles y las que más necesitan
    apoyo diagnóstico.
\end{enumerate}

En conclusión, reutilizar el modelo es técnicamente factible y recomendable desde una
perspectiva de eficiencia, pero éticamente inaceptable sin un proceso de validación que
confirme que el modelo es seguro para el nuevo dominio. La tecnología es transferible; la
confianza clínica debe ganarse desde cero.

% ===================== TASK 4 =====================
\newpage
\section*{\LARGE Task 4: Comentario en LinkedIn}

A continuación se muestra la captura de pantalla del comentario realizado en la publicación
de LinkedIn del curso, junto con las reacciones de los demás integrantes del grupo.

\begin{figure}[h!]
    \centering
    \includegraphics[width=0.85\textwidth]{comentario_linkedin.png}
    \caption{Captura de pantalla del comentario en la publicación de LinkedIn del curso.}
\end{figure}

% ===================== FIN =====================
\vfill
\begin{center}
    \small CC3182 -- Visión por Computadora $\cdot$ Universidad del Valle de Guatemala $\cdot$ 2026
\end{center}

\end{document}
